\documentclass[addpoints,spanish, 12pt,a4paper]{exam}
\pointpoints{punto}{puntos}
\hpword{Puntos:}
\vpword{Puntos:}
\htword{Total}
\vtword{Total}
\hsword{Resultado:}
\hqword{Ejercicio:}
\vqword{Ejercicio:}

\printanswers

\usepackage[utf8]{inputenc}
\usepackage[spanish]{babel}
\usepackage{eurosym}
%\usepackage[spanish,es-lcroman, es-tabla, es-noshorthands]{babel}


\usepackage[margin=1in]{geometry}
\usepackage{amsmath,amssymb}
\usepackage{multicol}

\usepackage{graphicx}
\graphicspath{{../img/}} 

\newcommand{\class}{Matemáticas 4º Aplicadas}
\newcommand{\examdate}{\today}
\newcommand{\examnum}{Enteros, fracciones y potencias}
\newcommand{\tipo}{A}


\newcommand{\timelimit}{50 minutos}

\renewcommand{\solutiontitle}{\noindent\textbf{Solución:}\enspace}

\pagestyle{head}
\firstpageheader{\includegraphics[width=0.2\columnwidth]{header_left}}{\textbf{Departamento de Matemáticas\linebreak \class}\linebreak \examnum}{\includegraphics[width=0.1\columnwidth]{header_right}}
\runningheader{\class}{\examnum}{Página \thepage\ of \numpages}
\runningheadrule


\begin{document}

\noindent
\begin{tabular*}{\textwidth}{l @{\extracolsep{\fill}} r @{\extracolsep{6pt}} }
\textbf{Nombre:} \makebox[3.5in]{\hrulefill} & \textbf{Fecha:}\makebox[1in]{\hrulefill} \\
 & \\
\textbf{Tiempo: \timelimit} & Tipo: \tipo 
\end{tabular*}
\rule[2ex]{\textwidth}{2pt}
Esta prueba tiene \numquestions\ ejercicios. La puntuación máxima es de \numpoints. 
La nota final de la prueba será la parte proporcional de la puntuación obtenida sobre la puntuación máxima. Para la evaluación de pendientes de 3ºESO o 2ºPMAR se tendrán en cuenta los apartados: 1.a, 2.a, 3.b y 4a,b,c

\begin{center}


\addpoints
 %\gradetable[h][questions]
	\pointtable[h][questions]
\end{center}

\noindent
\rule[2ex]{\textwidth}{2pt}

\begin{questions}

\question[3] Calcula:
\noaddpoints % to omit double points count
\begin{parts}
\part[1] $ -3 + 7 [-4 - (-12) : (-6) + 4 \cdot (-3)]$ 
\begin{solution}$= -3 + 7 [-4 - 2 - 12] = -3 + 7 \cdot (-18) = -3 - 126 = -129$ \end{solution}
\part[1] $14 : (-2) + (-5) : 5 - (-3) + 12$
\begin{solution}$=-7 - 1 + 3 + 12 = 7$ \end{solution}
\part[1] $7\cdot\left[12 + \left(-6 + 4 + 8\right)\right] - \left(-2\right) \cdot \left[5 - 3 \cdot \left(2 + 3 - 6\right)\right]$
\begin{solution}$= 7 [12 + 6] + 2 \cdot [5 - 3 \cdot (-1)] = 7 \cdot [18] + 2 \cdot [5 + 3] = 126 + 2 \cdot 8 = 126 + 16 = 142$ \end{solution}
\end{parts}
\addpoints

\question[3] Efectúa y simplifica:
\noaddpoints % to omit double points count
\begin{parts}
\part[1] $\dfrac{3}{2}-\dfrac{4}{5}:\dfrac{1}{2}+\dfrac{3}{4}\cdot\dfrac{1}{3}$ 
\begin{solution}$= \dfrac{3}{2}-\dfrac{8}{5}+\dfrac{1}{4}=\dfrac{30}{20}-\dfrac{32}{20}+\dfrac{5}{20}=\dfrac{3}{20}$ \end{solution}
\part[1] $\dfrac{1}{6}-\dfrac{5}{3}\left(\dfrac{4}{5}-\dfrac{1}{3}\right)-\dfrac{1}{2}:\dfrac{3}{4}$
\begin{solution}$=\dfrac{1}{6}-\dfrac{5}{3}\cdot\dfrac{7}{15}-\dfrac{2}{3}=\dfrac{1}{6}-\dfrac{7}{9}-\dfrac{2}{3}=\dfrac{3}{18}-\dfrac{14}{18}-\dfrac{12}{18}=-\dfrac{23}{18}$ \end{solution}
\part[1] $\left(\dfrac{3}{5}-\dfrac{2}{3}\right)^3\cdot\left[\left(\dfrac{1}{2}+\dfrac{1}{3}\right)^3:\left(\dfrac{5}{3}-1\right)^3\right]$
\begin{solution}$=\left(\dfrac{9-10}{15}\right)^3\cdot\left[\left(\dfrac{3+2}{6}\right)^3:\left(\dfrac{2}{3}\right)^3\right]=$ \\ 
$=\left(-\dfrac{1}{15}\right)^3\cdot\left[\left(\dfrac{5}{6}\right)^3:\left(\dfrac{2}{3}\right)^3\right]=\left(-\dfrac{1}{15}\right)^3\cdot\left(\dfrac{5}{4}\right)^3=\left(-\dfrac{1}{12}\right)^3=-\dfrac{1}{1728}$ \end{solution}
\end{parts}
\addpoints

\question[2] Simplifica utilizando las propiedades de las potencias:
\noaddpoints % to omit double points count
\begin{parts}
\part[1] $\dfrac{3^4\cdot 3 \cdot 9^2}{3^0\cdot 3 \cdot 27}$ 
\begin{solution}$=\dfrac{3^5\cdot  \left(3^2\right)^2}{3 \cdot \left(3^3\right)}=\dfrac{3^5\cdot 3^4}{3^4}=3^5=243$ \end{solution}
\part[1] $\dfrac{6^3\cdot 3^{-2}}{3^6\cdot 2^{-2}}$
\begin{solution}$=\dfrac{\left(2\cdot3\right)^3\cdot 2^{2}}{3^6\cdot 3^{2}}=\dfrac{\left(2^3\cdot3^3\right)\cdot 2^{2}}{3^6\cdot 3^{2}}=\dfrac{2^5\cdot3^3}{3^8}=\left(\dfrac{2}{3}\right)^5$ \end{solution}
\end{parts}
\addpoints

\question[3] Juan se gasta 2/3 del dinero en ropa y 1/4 del total en comida.:
\noaddpoints % to omit double points count
\begin{parts}
\part[1] ¿Cuál es la fracción gastada?
\begin{solution}$\frac{2}{3}+\frac{1}{4}=\frac{8+3}{12}=\frac{11}{12}$ \end{solution}
\part[1] ¿Qué fracción le queda por gastar?
\begin{solution}$1-\frac{11}{12}=\frac{1}{12}$ \end{solution}
\part[1] Si salió de casa con 180 \euro, ¿qué cantidad no se ha gastado?
\begin{solution} $\dfrac{1}{12} \ de \ 180=\dfrac{1\cdot180}{12}=15$ \euro \end{solution}
\end{parts}
\addpoints

\question[2] Un jardinero riega en un día 2/5 partes del jardín. ¿Cuántos días tardará en regar todo el jardín? ¿Cuánto ganará si cobra 50 \euro. por día?
\begin{solution} Días tardará: $1:\frac{2}{5}=\frac{5}{2}=2,5$ días  \\
Dinero que ganará: $50\cdot2,5=125$ \euro. \end{solution}


\question[2] Pablo gasta 2/5 del dinero que tenía en comprar fruta. Después, gasta 1/4 de lo que le queda en comprar leche. Sabiendo que le han sobrado 9 \euro. ¿Cuánto dinero tenía al principio?
\begin{solution} Fracción gastada: $\frac{2}{5}+\frac{1}{4}\cdot\frac{3}{5}=\frac{8}{20}+\frac{3}{20}=\frac{11}{20}$  \\
Fracción que le queda: $1-\frac{11}{20}=\frac{9}{20}$ \\
Dinero que tenía: Si $\frac{9}{20}$ del Total $= 9\to$ Total $=\dfrac{9\cdot20}{9}=20$ \euro. \end{solution}

\addpoints


\end{questions}

\end{document}
